\scanpage{72}
$\unsure{\underline{\mathbb W}|\underline W}^a(\unsure{\mathfrak X|\mathcal X}_i\cap\mathfrak X_j)$ est \inserted{donc} \inserted{à la fois \unsure{un}} ouvert \deleted{à la fois} \inserted{$X_{ij}$} dans $\underline{\mathbb W}^a(\mathfrak X_i)$ et \inserted{$X_{ji}$ dans} $\mathbb W^a(\mathfrak X_j)$ : cela permet de recoller les $\mathbb W^a(\mathfrak X_i)=X_i$. Il faut vérifier pour cela la \deleted{\unsure{$X_{ij}$ et $X_{jk}$, et|$X_{ijk}=X_{ikj}$}} commutativité \unsure{de|du}.
\diagram{\begin{tikzcd}[ampersand replacement=\&,column sep=large,row sep=small]
 X_{ij\unsure{h|k}} \arrow[r,"{\unsure{\varphi|\phi}_{ji}\vert}"] \arrow[rr,bend left=25,"{\varphi_{\unsure{h|k}i}\vert}"] \& X_{jih} \arrow[r,"{\varphi_{kj}\vert}"] \& X_{kij}\\
 \parallel \& \parallel \& \parallel\\
 X_{ij}\cap X_{ih} \& X_{ji}\cap X_{jh} \& X_{ki}\cap X_{kj}
\end{tikzcd}}

Cela tient au fait que l'on aura $X_{ijh}\simeq\underline{\mathbb W}^a(\mathfrak X_i\cap\mathfrak X_j\cap\mathfrak X_k)$ et que \deleted{les} \unsure{aplication|application|applications} $\varphi_{ji}$ (p. ex) s'identifie à l'identité.

On a ainsi défini un $X=\bigcup X_i$, $X_i\simeq\underline{\mathbb W}^a(\mathfrak X_i)$. On définit
\[
 \mathbb W(X)\longrightarrow\mathfrak X
\]
en notant que $\mathbb W(X)=\bigcup_i\mathbb W(X_i)$ \deleted{\unsure{et|$\simeq$}} et qu'on a des \deleted{\unsure{hom}} morphismes
\[
 \underline{\mathbb W}(X_i)\longrightarrow\mathfrak X_i\subset\mathfrak X
\]
par définition de $X_i=\underline{\mathbb W}^a(\mathfrak X_i)$. La
% archive sheet number: 71
